% =====================================================================
%  pds-comun.tex — Paquetes, colores y estilos de código comunes a
%  transparencias, apuntes y hojas de ejercicios. Compilar con XeLaTeX.
%  No cargar directamente: lo cargan pds-beamer.tex y pds-apuntes.tex.
% =====================================================================

\usepackage{polyglossia}
\setmainlanguage{spanish}
\usepackage{mathtools}
\usepackage{graphicx}
\usepackage{booktabs}
\usepackage{multirow}
\usepackage{array}
\usepackage{xcolor}
\usepackage{tikz}
\usetikzlibrary{positioning,arrows.meta,calc,shapes.geometric,decorations.pathreplacing}
\usepackage{listings}
\usepackage{siunitx}
\sisetup{output-decimal-marker={,}, group-separator={\,}, per-mode=symbol}

% Notas de erratas de versiones anteriores: no se imprimen salvo que se
% defina \ERRATAS antes de cargar el estilo (make ERRATAS=1)
\newif\iferratas
\ifdefined\ERRATAS \erratastrue \fi

% Ruta común (logo y figuras compartidas)
\providecommand{\rutacomun}{../../comun/}
\input{\rutacomun datos.tex}

% ---------------------------------------------------------------------
%  Paleta (la misma que usan las figuras generadas con Python,
%  ver comun/pds_estilo.py)
% ---------------------------------------------------------------------
\definecolor{pdsazul}{HTML}{2B3A8C}
\definecolor{pdsrojo}{HTML}{C0392B}
\definecolor{pdsverde}{HTML}{1E7B4F}
\definecolor{pdsnaranja}{HTML}{D9822B}
\definecolor{pdsgris}{HTML}{5A5A66}
\definecolor{pdsfondo}{HTML}{F3F4FA}
\definecolor{pdscodigo}{HTML}{F7F7F2}

% ---------------------------------------------------------------------
%  Diagramas de bloques (TikZ)
% ---------------------------------------------------------------------
\tikzset{
  bloque/.style={rectangle, rounded corners=2pt, draw=pdsazul!80, fill=pdsfondo,
                 very thick, minimum height=8mm, minimum width=14mm, align=center},
  sumador/.style={circle, draw=pdsazul!80, fill=pdsfondo, thick, inner sep=1.5pt},
  flecha/.style={-{Stealth[length=2.2mm]}, thick, pdsgris},
}

% ---------------------------------------------------------------------
%  Código: Python y Matlab en paralelo
% ---------------------------------------------------------------------
\newcommand{\pdscodesize}{\footnotesize}
\lstdefinestyle{pds}{
  basicstyle=\ttfamily\pdscodesize,
  keywordstyle=\color{pdsazul}\bfseries,
  commentstyle=\color{pdsverde}\itshape,
  stringstyle=\color{pdsrojo},
  numbers=none,
  backgroundcolor=\color{pdscodigo},
  frame=single, rulecolor=\color{pdsgris!40}, framerule=0.4pt,
  xleftmargin=2pt, xrightmargin=2pt, aboveskip=2pt, belowskip=2pt,
  columns=fullflexible, keepspaces=true,
  showstringspaces=false, breaklines=true, tabsize=4,
  upquote=true,
  % Con XeLaTeX, listings desordena los caracteres no ASCII junto a
  % dígitos («2π» sale «π2»): se sustituyen por su versión matemática.
  literate={π}{{$\pi$}}1 {δ}{{$\delta$}}1 {ω}{{$\omega$}}1 {Ω}{{$\Omega$}}1
           {≈}{{$\approx$}}1 {µ}{{$\mu$}}1 {·}{{$\cdot$}}1 {α}{{$\alpha$}}1
           {á}{{á}}1 {é}{{é}}1 {í}{{í}}1 {ó}{{ó}}1 {ú}{{ú}}1 {ñ}{{ñ}}1 {¿}{{¿}}1,
}
\lstdefinestyle{python}{style=pds, language=Python,
  morekeywords={as,True,False,None}}
\lstdefinestyle{matlab}{style=pds, language=Matlab,
  morekeywords={function,end,true,false}}

% \codigo{<ruta sin extensión>} muestra el mismo ejemplo en Python
% y en Matlab, uno al lado del otro. Los ficheros viven en
% ../codigo/python/<nombre>.py y ../codigo/matlab/<nombre>.m y se
% ejecutan en las pruebas (make prueba-codigo).
\newcommand{\rutacodigo}{../codigo/}
\newcommand{\codigoetiqueta}[1]{{\sffamily\scriptsize\bfseries\color{pdsgris}#1}}
\newcommand{\codigo}[2][]{%
  \par\noindent
  \begin{minipage}[t]{0.495\linewidth}
    \codigoetiqueta{Python}\par\vspace{1pt}
    \lstinputlisting[style=python,#1]{\rutacodigo python/#2.py}
  \end{minipage}\hfill
  \begin{minipage}[t]{0.495\linewidth}
    \codigoetiqueta{Matlab}\par\vspace{1pt}
    \lstinputlisting[style=matlab,#1]{\rutacodigo matlab/#2.m}
  \end{minipage}\par}

\input{\rutacomun notacion.tex}
